\documentclass[11pt]{article}
\usepackage[a4paper,margin=24mm]{geometry}
\usepackage[T1]{fontenc}
\usepackage{amsmath,amssymb,amsthm,mathtools}
\usepackage[hidelinks,hypertexnames=false]{hyperref}
\setlength{\emergencystretch}{2em}
\newcommand{\C}{\mathbb C}\newcommand{\R}{\mathbb R}
\newcommand{\tr}{\operatorname{tr}}\newcommand{\dd}{\dagger_G}
\newcommand{\eps}{\epsilon}\newcommand{\spec}{\operatorname{spec}}
\newtheorem{theorem}{Theorem}[section]
\newtheorem{lemma}[theorem]{Lemma}
\title{The finite native heat action\\
Convergent arithmetic secular coefficients and retained metric control}
\author{}\date{20 September 2026}
\begin{document}\maketitle

\paragraph{Scope and human provenance.}
This note extends the original finite arithmetic path of
\cite{NativeAction}%
% BEGIN VERIFIED PUBLIC PROOF LINK
~\href{https://github.com/KokunoYumeto/zeta-function-research-reader/blob/5992ad50c0f2a06921f1fcaded7b4e38097c0739/workbenches/splitzero-tandem/continuations/20260920-original-kernel-mixed-determinants/providers/NATIVE_SPECTRAL_ACTION.tex\#L31}{[proof SA1]} \href{https://github.com/KokunoYumeto/zeta-function-research-reader/blob/5992ad50c0f2a06921f1fcaded7b4e38097c0739/workbenches/splitzero-tandem/continuations/20260920-original-kernel-mixed-determinants/providers/SOURCE_FUNCTION_CLASS_CHECK.tex\#L24}{[proof FC1]} \href{https://github.com/KokunoYumeto/zeta-function-research-reader/blob/5992ad50c0f2a06921f1fcaded7b4e38097c0739/workbenches/splitzero-tandem/continuations/20260920-original-kernel-mixed-determinants/providers/RANK_ONE_METRIC_CURVATURE.tex\#L35}{[proof MC1]} \href{https://github.com/KokunoYumeto/zeta-function-research-reader/blob/5992ad50c0f2a06921f1fcaded7b4e38097c0739/workbenches/splitzero-tandem/continuations/20260920-original-kernel-mixed-determinants/providers/INDEPENDENT_ENDPOINT_CHECK.tex\#L28}{[proof EC1]}%
% END VERIFIED PUBLIC PROOF LINK
{}, SA1--SA43, from polynomial tests to the
genuine heat function \(F_s(u)=e^{-su^2}\).
The Gaussian multiplicative Fourier transform is exactly the
one used by Alain Connes \cite{ConnesHeat}, equation
\texttt{ft} in the original author source. Its finite matrix
application is proved below; Connes's infinite selfadjoint
heat theorem, labelled \texttt{heatkernelintro}, explicitly
assumes RH and is not invoked here.
The finite divided-difference method is discussed by Connes
and van Suijlekom \cite{CV}; the receiving polynomial
calculation and its exact cyclic factor are proved in
\cite{NativeAction}%
% BEGIN VERIFIED PUBLIC PROOF LINK
~\href{https://github.com/KokunoYumeto/zeta-function-research-reader/blob/5992ad50c0f2a06921f1fcaded7b4e38097c0739/workbenches/splitzero-tandem/continuations/20260920-original-kernel-mixed-determinants/providers/NATIVE_SPECTRAL_ACTION.tex\#L31}{[proof SA1]} \href{https://github.com/KokunoYumeto/zeta-function-research-reader/blob/5992ad50c0f2a06921f1fcaded7b4e38097c0739/workbenches/splitzero-tandem/continuations/20260920-original-kernel-mixed-determinants/providers/SOURCE_FUNCTION_CLASS_CHECK.tex\#L24}{[proof FC1]} \href{https://github.com/KokunoYumeto/zeta-function-research-reader/blob/5992ad50c0f2a06921f1fcaded7b4e38097c0739/workbenches/splitzero-tandem/continuations/20260920-original-kernel-mixed-determinants/providers/RANK_ONE_METRIC_CURVATURE.tex\#L35}{[proof MC1]} \href{https://github.com/KokunoYumeto/zeta-function-research-reader/blob/5992ad50c0f2a06921f1fcaded7b4e38097c0739/workbenches/splitzero-tandem/continuations/20260920-original-kernel-mixed-determinants/providers/INDEPENDENT_ENDPOINT_CHECK.tex\#L28}{[proof EC1]}%
% END VERIFIED PUBLIC PROOF LINK
{}, SA8--SA11. All heat formulas and error
bounds required here have complete proofs below.
This is result SZ-20260920-043. Every trace in this note is
finite; no infinite spectral realization or zero-free
conclusion is asserted.

\section{The original arithmetic space, without changing its metric}

Keep the nonempty finite packet \(Z\) of actual zeros with
their full orders \(m_\rho\), invariant under conjugation
and \(\rho\mapsto1-\rho\). The source data are
\begin{equation}\tag{NH1}\label{NH1}
\begin{gathered}
h(v)=\prod_{\rho\in Z}(v-\rho)^{m_\rho},\qquad
g(v)=2\xi(v)=v(v-1)\pi^{-v/2}\Gamma(v/2)\zeta(v),\\
\mathcal M F_h=g/h,\qquad
w_h(u)=\frac{|(g/h)(1/2+iu)|^2}{2\pi},\qquad
m(u)=w_h^{*k}(u),\qquad c=k/2,\quad S=c+iu .
\end{gathered}
\end{equation}
Use the actual positive even measure \(m(u)\,du\) with
its full mass
\(\mathfrak m=(\int_\R w_h)^k\). Its positivity almost
everywhere and finite moments are the original source
properties retained in \cite{NativeAction}%
% BEGIN VERIFIED PUBLIC PROOF LINK
~\href{https://github.com/KokunoYumeto/zeta-function-research-reader/blob/5992ad50c0f2a06921f1fcaded7b4e38097c0739/workbenches/splitzero-tandem/continuations/20260920-original-kernel-mixed-determinants/providers/NATIVE_SPECTRAL_ACTION.tex\#L31}{[proof SA1]} \href{https://github.com/KokunoYumeto/zeta-function-research-reader/blob/5992ad50c0f2a06921f1fcaded7b4e38097c0739/workbenches/splitzero-tandem/continuations/20260920-original-kernel-mixed-determinants/providers/SOURCE_FUNCTION_CLASS_CHECK.tex\#L24}{[proof FC1]} \href{https://github.com/KokunoYumeto/zeta-function-research-reader/blob/5992ad50c0f2a06921f1fcaded7b4e38097c0739/workbenches/splitzero-tandem/continuations/20260920-original-kernel-mixed-determinants/providers/RANK_ONE_METRIC_CURVATURE.tex\#L35}{[proof MC1]} \href{https://github.com/KokunoYumeto/zeta-function-research-reader/blob/5992ad50c0f2a06921f1fcaded7b4e38097c0739/workbenches/splitzero-tandem/continuations/20260920-original-kernel-mixed-determinants/providers/INDEPENDENT_ENDPOINT_CHECK.tex\#L28}{[proof EC1]}%
% END VERIFIED PUBLIC PROOF LINK
{}, SA1--SA2.
Let \(Q_j\) be its real monic orthogonal polynomials,
\(\omega_j=\int Q_j^2m\), and
\(p_j(S)=i^jQ_j((S-c)/i)\). Thus
\(\omega_0=\mathfrak m\), not \(1\).

The cyclic arithmetic algebra and its complete relation are
\begin{equation}\tag{NH2}\label{NH2}
C=\C[S]/(\chi),\qquad
\chi(S)=\prod_\kappa(S-\kappa)^{e_\kappa},\qquad
e_\kappa=\max_{\sum_{j=1}^k\rho_j=\kappa}
\left(1+\sum_{j=1}^k(m_{\rho_j}-1)\right),\qquad q=\deg\chi .
\end{equation}
Here is the calculation retaining collisions. In the local
tensor ring with coordinates \(z_j^{m_{\rho_j}}=0\),
the highest nonzero power of \(\sum z_j\) has exponent
\(\sum(m_{\rho_j}-1)\). Its coefficient on
\(\prod z_j^{m_{\rho_j}-1}\) is
\((\sum(m_{\rho_j}-1))!/\prod(m_{\rho_j}-1)!\ne0\).
The next power is zero. The least polynomial annihilating
all local tensor components therefore has exactly the
maximum exponent at each repeated sum in \eqref{NH2}.

Fix the original degree \(L\ge q-1\), and let \(b_j=[p_j]_\chi\).
In the ordered remainder basis \(1,S,\ldots,S^{q-1}\), set
\begin{equation}\tag{NH3}\label{NH3}
\begin{gathered}
E=[i^{-j}b_j/\sqrt{\omega_j}]_{j=0}^L,\quad
G=(EE^*)^{-1},\quad R=E^*G,\quad ER=I,\quad R^*R=G,\\
M=(M_S-cI)/i,\qquad
A=EJ_LR=M+Q,\qquad
Q=r\ell,\quad r=\frac{i}{\omega_L}b_{L+1},\quad
\ell=b_L^*G .
\end{gathered}
\end{equation}
Here \(J_L\) is compression of multiplication by \(u\)
in the orthonormal polynomial basis
\(Q_j/\sqrt{\omega_j}\).
The first \(q\) monic columns of \(E\) are independent,
so \(EE^*>0\). A lift of \(v\in C\) has form \(Rv+z\),
\(Ez=0\), and
\(\|Rv+z\|^2=v^*Gv+\|z\|^2\), since
\((Rv)^*z=v^*GEz=0\).
Thus this is precisely the attained quotient metric.

For completeness, the only extra column in full multiplication
is \(\sqrt{\omega_{L+1}/\omega_L}\,e_{L+1}e_L^*\).
Reduction and lifting turn it into
\[
\sqrt{\omega_{L+1}/\omega_L}\,
\frac{i^{-(L+1)}b_{L+1}}{\sqrt{\omega_{L+1}}}
\frac{i^Lb_L^*G}{\sqrt{\omega_L}}
=-\frac{i}{\omega_L}b_{L+1}b_L^*G .
\]
Removing that column proves the sign in \eqref{NH3}.
Since \(GA=R^*J_LR\) is Hermitian,
\(A^\dd=A\), where \(X^\dd=G^{-1}X^*G\).

Parity gives \(Q_j(-u)=(-1)^jQ_j(u)\).
The actual involution \(\mathcal P[P(S)]=[P(k-S)]\)
satisfies
\begin{equation}\tag{NH4}\label{NH4}
\mathcal P^*G\mathcal P=G,\quad
\mathcal P A\mathcal P=-A,\quad
\mathcal P Q\mathcal P=-Q,\quad
\ell r=0,\quad Q^2=0 .
\end{equation}
Indeed \(\mathcal P E=E\operatorname{diag}((-1)^j)\);
insert this into \eqref{NH3}. Orthogonal polynomials of
opposite parity are orthogonal in the quotient metric,
because the displayed involution is an isometry.
This proves \(b_L^*Gb_{L+1}=0\), hence the last two
identities. Put
\begin{equation}\tag{NH5}\label{NH5}
T(t)=A-tQ,\quad T(1)=M,\quad
a=\|A\|_G,\quad
\eps=\|Q\|_G
=\frac{\sqrt{(b_L^*Gb_L)(b_{L+1}^*Gb_{L+1})}}{\omega_L},
\quad \gamma_j=\ell A^jr .
\end{equation}
The rank-one norm formula follows by writing
\(Qx=r\langle G^{-1}\ell^*,x\rangle_G\) and attaining
Cauchy--Schwarz on \(G^{-1}\ell^*\). It also gives
\(\tr Q^\dd Q=\eps^2\) and
\(|\gamma_j|\le\eps a^j\).
Parity in \eqref{NH4} proves \(\gamma_{2j}=0\).
Complex \(t\) is allowed in the holomorphic formulas.
Statements involving adjoints of \(T(t)\) use real \(t\).

\section{A finite heat function with the exact Fourier constants}

For \(s>0\), the Gaussian formula is
\begin{equation}\tag{NH6}\label{NH6}
F_s(z)=e^{-sz^2}
=\frac1{\sqrt{4\pi s}}\int_\R e^{-v^2/(4s)}e^{ivz}\,dv,
\qquad z\in\C .
\end{equation}
To prove it, the integral of \(e^{-x^2}\) is \(\sqrt\pi\)
by squaring it and using polar coordinates. Differentiating
the right side of \eqref{NH6} in \(z\) and integrating
by parts in \(v\) gives the differential equation
\(J'(z)=-2szJ(z)\), with \(J(0)=1\).
On every compact set in \(z\), a polynomial times
\(e^{-v^2/(4s)+|v||z|}\) is integrable, justifying both
operations and analyticity. This proves \eqref{NH6}
including nonreal \(z\). Under \(x=e^v\), its Gaussian
is exactly Connes's \texttt{ft} \cite{ConnesHeat};
no factor \(2\), \(\pi\), or square root has been removed.

In the fixed quotient norm, matrix power series give
\begin{equation}\tag{NH7}\label{NH7}
\Phi_s(t):=\tr e^{-sT(t)^2}
=\sum_{m=0}^\infty\frac{(-s)^m}{m!}\tr T(t)^{2m}
=\frac1{\sqrt{4\pi s}}\int_\R
e^{-v^2/(4s)}\tr e^{ivT(t)}\,dv .
\end{equation}
All three expressions are equal. For the first series,
the norm of its \(m\)-th matrix term is at most
 \(s^m\|T(t)\|_G^{2m}/m!\).
For the Fourier expression expand \(e^{ivT(t)}\) and
use the dominating exponential
\(e^{|v|\|T(t)\|_G-v^2/(4s)}\).
Odd Gaussian moments vanish and the \(2m\)-th moment is
\((2m)!s^m/m!\), by integration by parts. This proves
the equality term by term. On compact sets of \(s,t\),
the first series converges locally uniformly, including
complex \(s\). Thus \(\Phi\) is entire in both variables.
At any finite degree \(D\ge0\), with \(b=\|T(t)\|_G\),
\begin{equation}\tag{NH8}\label{NH8}
\left|\Phi_s(t)-\sum_{m=0}^D
\frac{(-s)^m}{m!}\tr T(t)^{2m}\right|
\le q e^{|s|b^2}\frac{(|s|b^2)^{D+1}}{(D+1)!}.
\end{equation}
Indeed \(|\tr X|\le q\|X\|_G\), by diagonal entries in
a \(G\)-orthonormal basis, and
\(\sum_{m>D}x^m/m!\le e^x x^{D+1}/(D+1)!\);
the latter follows from
\((D+1+j)!\ge(D+1)!\,j!\).

\section{Full secular determinant and every Jordan endpoint}

Retain the full polynomials, not a reduced rational factor:
\begin{equation}\tag{NH9}\label{NH9}
\begin{gathered}
d(z)=\det(zI-A),\quad
\chi_u(z)=i^{-q}\chi(c+iz)=\det(zI-M),\\
n(z)=\ell\,\operatorname{adj}(zI-A)r=\chi_u(z)-d(z),\\
\kappa(z)=\ell(zI-A)^{-1}r=\frac{n(z)}{d(z)},\\
p_t(z)=\det(zI-T(t))=d(z)+tn(z)=(1-t)d(z)+t\chi_u(z).
\end{gathered}
\end{equation}
To prove the determinant identity away from the spectrum
of \(A\), factor \(zI-T=(zI-A)(I+t(zI-A)^{-1}r\ell)\).
The rank-one determinant is \(1+t\ell(zI-A)^{-1}r\),
as can be checked in a basis beginning with that column.
Multiplication by \(d\) makes it a polynomial identity
valid everywhere. At \(t=1\) it proves the numerator
identity. The companion matrix in the ordered powers of
\(S\) gives the formula for \(\chi_u\), with its full
factor \(i^{-q}\).

For every circle \(\mathcal C_R:|z|=R\) with
\(R>a+|t|\eps\), the exact heat formula is
\begin{equation}\tag{NH10}\label{NH10}
\Phi_s(t)=\frac1{2\pi i}\oint_{\mathcal C_R}
F_s(z)\frac{p_t'(z)}{p_t(z)}\,dz .
\end{equation}
Here and below the circle is oriented counterclockwise.
For proof expand \((zI-T)^{-1}=\sum_{j\ge0}T^jz^{-j-1}\)
uniformly on this circle. Multiplication by the uniformly
convergent entire series for \(F_s\) and integration gives
\(F_s(T)\). Triangularize the finite matrix \(T\);
the diagonal of its inverse is \(1/(z-\lambda)\).
Consequently its trace is \(p_t'/p_t\), with every
algebraic multiplicity, proving \eqref{NH10}.
This proof requires no diagonalizability.

The rank-one resolvent, checked by multiplication, gives
\begin{equation}\tag{NH11}\label{NH11}
\begin{gathered}
(zI-T(t))^{-1}
=(zI-A)^{-1}
-\frac{t(zI-A)^{-1}r\ell(zI-A)^{-1}}{1+t\kappa(z)},\\
\kappa_t(z)=\ell(zI-T(t))^{-1}r
=\frac{\kappa(z)}{1+t\kappa(z)}=\frac{n(z)}{p_t(z)},\\
\Phi_s^{(j)}(t)
=\frac{(-1)^j(j-1)!}{2\pi i}
\oint_{\mathcal C_R}F_s'(z)\kappa_t(z)^j\,dz,\qquad j\ge1.
\end{gathered}
\end{equation}
To prove the last formula, first use the trace of the
first identity to obtain
\(\tr(zI-T)^{-1}-\tr(zI-A)^{-1}
=\partial_z\log(1+t\kappa(z))\).
The logarithm on the circle is its convergent power series,
since \(|t\kappa(z)|\le|t|\eps/(R-a)<1\).
Integrate by parts around the circle and differentiate
in \(t\). This produces exactly
\((-1)^j(j-1)!\kappa_t^j\).
In a neighborhood of any fixed \(t\), choose a larger
circle with the same inequality, which justifies every
derivative. Formula \eqref{NH11} therefore holds at \(t=1\)
even if the original multiplication matrix has Jordan blocks.

Explicitly, on a generalized eigenspace with eigenvalue
\(\lambda\) and nilpotent part \(N_\lambda\) of exponent
\(d_\lambda\), the operator and its trace are
\begin{equation}\tag{NH12}\label{NH12}
F_s(T)|_\lambda
=\sum_{j=0}^{d_\lambda-1}
\frac{F_s^{(j)}(\lambda)}{j!}N_\lambda^j,\qquad
\Phi_s(t)=\sum_\lambda m_\lambda e^{-s\lambda^2}.
\end{equation}
The operator formula follows by substituting
\(\lambda I+N_\lambda\) into the entire Taylor series;
all terms \(j\ge d_\lambda\) are zero.
A basis adapted to the kernels of powers of \(N_\lambda\)
makes its positive powers strictly triangular, with trace
zero. Its constant term has trace \(m_\lambda F_s(\lambda)\).
Thus nilpotent directions have not been deleted from the
operator, and their dimensions remain in its trace.

\section{The parity-improved, absolutely convergent heat coefficients}

The original parity gives the convergent Laurent expansion
\begin{equation}\tag{NH13}\label{NH13}
\kappa(z)=\sum_{j=0}^\infty
\gamma_{2j+1}z^{-2j-2},\qquad |z|>a .
\end{equation}
This is the geometric resolvent expansion with
\(\gamma_{2j}=0\), proved in \eqref{NH4}--\eqref{NH5}.
For \(n\ge1\) and \(r\ge0\), define the finite ordered sum
\begin{equation}\tag{NH14}\label{NH14}
H_{n,r}=
\sum_{\substack{j_1,\ldots,j_n\ge0\\j_1+\cdots+j_n=r}}
\prod_{\nu=1}^n\gamma_{2j_\nu+1}.
\end{equation}
Every ordered composition is retained.

\begin{theorem}[Exact coefficients and explicit factorial bounds]
For all complex \(s,t\),
\begin{equation}\tag{NH15}\label{NH15}
\Phi_s(t)=\Phi_s(0)+\sum_{n=1}^\infty t^n C_n(s),\qquad
C_n(s)=2\sum_{r=0}^\infty
\frac{(-1)^r s^{n+r}}{n(n+r-1)!}\,H_{n,r}.
\end{equation}
For \(a=\|A\|_G\) and the original \(\eps=\|Q\|_G\),
\begin{equation}\tag{NH16}\label{NH16}
|C_n(s)|\le
2e^{|s|a^2}\frac{(|s|a\eps)^n}{n!},\qquad n\ge1.
\end{equation}
If \(a=0\) or \(\eps=0\), all the \(C_n\) vanish.
These formulas do not divide by either number.
\end{theorem}
\begin{proof}
Expand the logarithm in the proof of \eqref{NH11}
uniformly on its circle. The coefficient of \(t^n\) is
\[
\frac{(-1)^n}{n}\frac1{2\pi i}
\oint F_s'(z)\kappa(z)^n\,dz .
\]
The entire derivative is
\(F_s'(z)=\sum_{m\ge1}2m(-s)^m z^{2m-1}/m!\).
The \(n\)-th power of \eqref{NH13} is
\(\sum_{r\ge0}H_{n,r}z^{-2n-2r}\).
Both series converge absolutely and uniformly on the
circle, so integration selects \(m=n+r\).
The resulting coefficient is precisely \eqref{NH15}.
There are \(\binom{n+r-1}{n-1}\) ordered \(n\)-tuples
of nonnegative integers summing to \(r\): place \(n-1\)
separators among \(n+r-1\) positions. By \eqref{NH5},
\[
|H_{n,r}|\le
\binom{n+r-1}{n-1}\eps^n a^{n+2r}.
\]
Multiply this by \(2|s|^{n+r}/(n(n+r-1)!)\).
The factorials give exactly
\[
\frac{2(|s|a\eps)^n}{n!}\frac{(|s|a^2)^r}{r!}.
\]
Summing over \(r\) proves \eqref{NH16}.
It also proves absolute local uniform convergence of the
double series in \(s,t\). The initial coefficient identity,
valid in a neighborhood of \(t=0\), consequently extends
to every complex \(t\). When \(a=0\), the bound
\(|\gamma_{2j+1}|\le\eps a^{2j+1}\) makes every factor
zero, proving the stated degenerate case directly.
\end{proof}

Write \(x=|s|a^2\), \(y=|s||t|a\eps\).
For every \(N\ge0\), the actual perturbation remainder is
\begin{equation}\tag{NH17}\label{NH17}
\left|\Phi_s(t)-\Phi_s(0)-\sum_{n=1}^N t^n C_n(s)\right|
\le 2e^x\left(e^y-\sum_{n=0}^N\frac{y^n}{n!}\right)
\le2e^{x+y}\frac{y^{N+1}}{(N+1)!}.
\end{equation}
The first inequality sums \eqref{NH16}; the second
uses the factorial tail proof following \eqref{NH8}.
In particular
\(|\Phi_s(t)-\Phi_s(0)|\le2e^x(e^y-1)\).
This error has no factor \(q\); it is a consequence of
the actual rank-one relation and parity, not of discarding
the \(q\) original dimensions from \(\Phi_s(0)\).

There are explicit inner-series errors as well. Set
\[
C_{n,R}(s)=2\sum_{r=0}^{R}
\frac{(-1)^rs^{n+r}}{n(n+r-1)!}H_{n,r}.
\]
Then
\begin{equation}\tag{NH18}\label{NH18}
|C_n(s)-C_{n,R}(s)|
\le\frac{2(|s|a\eps)^n}{n!}
e^x\frac{x^{R+1}}{(R+1)!}.
\end{equation}
Thus evaluating only \(1\le n\le N,0\le r\le R\)
has total certified error bounded by the sum of
\eqref{NH17} and
\(2e^x x^{R+1}(e^y-1)/(R+1)!\).
The proof is the same exact product bound used above;
no unknown higher moment is replaced by zero.

The leading coefficient and its retained correction are
\begin{equation}\tag{NH19}\label{NH19}
C_n(s)=\frac{2(s\gamma_1)^n}{n!}
+\left(C_n(s)-\frac{2(s\gamma_1)^n}{n!}\right),\qquad
\left|C_n(s)-\frac{2(s\gamma_1)^n}{n!}\right|
\le\frac{2(|s|a\eps)^n}{n!}(e^x-1).
\end{equation}
For example,
\begin{equation}\tag{NH20}\label{NH20}
C_1(s)=2s\gamma_1-2s^2\gamma_3+\cdots,\qquad
C_2(s)=s^2\gamma_1^2-s^3\gamma_1\gamma_3+\cdots .
\end{equation}
The displayed \(s^3\) coefficient in \(C_2\) has two
ordered compositions and the denominator \(2(2!)\);
these give the factor \(1\), not \(1/2\).
Each omitted tail is governed by \eqref{NH18}.

One may instead truncate the heat degree \(m=n+r\).
Let \(\Delta_{s,D}(t)\) be the terms of
\(\Phi_s(t)-\Phi_s(0)\) with \(n+r\le D\).
For \(D\ge0\),
\begin{equation}\tag{NH21}\label{NH21}
|\Phi_s(t)-\Phi_s(0)-\Delta_{s,D}(t)|
\le 2y e^{x+y}\frac{(x+y)^D}{D!}.
\end{equation}
Indeed the absolute bound at total degree \(m\) is
\(2((x+y)^m-x^m)/m!\).
For nonnegative \(x,y\),
\((x+y)^m-x^m\le my(x+y)^{m-1}\);
sum \(m\ge D+1\) and apply the exponential tail estimate.
The bound vanishes with the actual perturbation \(tQ\).

\section{A separate Gaussian derivative bound}

The following second bound can be used when it is smaller.
For \(s>0,n\ge1\),
\begin{equation}\tag{NH22}\label{NH22}
|C_n(s)|\le
\frac{(\sqrt{s}\eps)^n}{\Gamma(n/2+1)} .
\end{equation}
Here is a direct proof with every rank and constant retained.
The coefficient of \(t^n\) in \(e^{iv(A-tQ)}\) is
\[
(-iv)^n\int_{\substack{\alpha_0,\ldots,\alpha_n\ge0\\
                       \sum\alpha_j=1}}
e^{iv\alpha_0A}Qe^{iv\alpha_1A}\cdots
Qe^{iv\alpha_nA}\,d\alpha .
\]
To verify the formula without an operator expansion theorem,
expand the \(n+1\) exponentials. Integration of
\(\prod\alpha_j^{m_j}\) over the displayed simplex equals
\(\prod m_j!/(m_0+\cdots+m_n+n)!\), by successive
one-variable beta integrals, each proved by integration
by parts. This is exactly the coefficient of each
noncommutative word in the matrix exponential series.
The simplex volume is \(1/n!\).
For real \(v\), all factors \(e^{iv\alpha A}\) have
norm \(1\) in \(G\). A rank-one operator \(r\ell\)
has trace norm \(\eps\), and
\(|\tr XY|\le\|X\|_1\|Y\|\); this follows by writing
the rank-one \(X=uv^*\) in orthonormal coordinates and
using Cauchy--Schwarz. Repeated use gives an absolute
trace bound \(|v|^n\eps^n/n!\).
Apply \eqref{NH6} and compute its absolute moment:
\[
\frac1{\sqrt{4\pi s}}\int_\R
e^{-v^2/(4s)}|v|^n\,dv
=\frac{(4s)^{n/2}\Gamma((n+1)/2)}{\sqrt\pi}.
\]
The identity
\(\Gamma((n+1)/2)/(\sqrt\pi n!)
=1/(2^n\Gamma(n/2+1))\) follows at integers by
\(\Gamma(1/2)=\sqrt\pi\) and the recurrence
\(\Gamma(z+1)=z\Gamma(z)\), or by induction on \(n\)
separately for its two parities.
This proves \eqref{NH22}. In particular the better of
\eqref{NH16} and \eqref{NH22} is available without
altering the packet, its mass, or its metric.

The repeated divided differences have their exact finite
spectral form too. If
\(A=\sum_{\lambda\in\Lambda}\lambda P_\lambda\) is its
complete \(G\)-selfadjoint spectral decomposition and
\(w_\lambda=\ell P_\lambda r\), then
\begin{equation}\tag{NH23}\label{NH23}
C_n(s)=\frac{(-1)^n}{n}
\sum_{\lambda_1,\ldots,\lambda_n\in\Lambda}
F_s'[\lambda_1,\ldots,\lambda_n]
\prod_{\nu=1}^n w_{\lambda_\nu}.
\end{equation}
Define the repeated divided difference by the circle
integral of \(F_s'(z)/\prod(z-\lambda_\nu)\).
At distinct nodes partial fractions give the usual
recursive formula; continuity extends it to repetitions,
where \(f[\lambda,\ldots,\lambda]=f^{(n-1)}(\lambda)/(n-1)!\).
Insert
\(\kappa(z)=\sum_\lambda w_\lambda/(z-\lambda)\)
in the coefficient integral above. This proves
\eqref{NH23}, with all repeated projectors retained.
Multiplication by \(n!\) gives the derivative factor
\((-1)^n(n-1)!\), not \((-1)^nn!\).

\section{Original scaling coordinate and the complete physical map}

Set \(B(t)=cI+iT(t)\). The original heat test in \(S\) is
\begin{equation}\tag{NH24}\label{NH24}
f_s(S)=e^{s(S-c)^2},\qquad
F_s(u)=f_s(c+iu),\qquad
\det(zI-B(t))=i^q p_t((z-c)/i).
\end{equation}
The determinant identity follows from
\(zI-B=i(((z-c)/i)I-T)\).
The positive sign in \(f_s\) comes from \(i^2=-1\).
The entire derivative and divided differences satisfy
\(F_s^{(j)}(u)=i^j f_s^{(j)}(c+iu)\) and
\[
F_s'[\lambda_1,\ldots,\lambda_n]
=i^n f_s'[c+i\lambda_1,\ldots,c+i\lambda_n].
\]
For distinct nodes this follows successively from the
divided-difference recursion; continuity proves every
coincidence. Hence the native scaling trace derivative
retains the factor \((-i)^n(n-1)!\) in its cyclic formula.
At the original endpoint,
\begin{equation}\tag{NH25}\label{NH25}
\Phi_s(1)=\tr f_s(M_S)
=\sum_\kappa e_\kappa e^{s(\kappa-c)^2}.
\end{equation}
This expression includes all collision exponents from
\eqref{NH2}, not just the distinct sums.

Let \(E_h=\C[v]/(h)\) and
\(\upsilon=j_h(g/h)\). Retain
\begin{equation}\tag{NH26}\label{NH26}
\eta:C\longrightarrow E_h^{\otimes k},\qquad
\eta[P]=\upsilon^{\otimes k}P(v_1+\cdots+v_k),\qquad
\langle\eta x,\eta y\rangle=x^*Gy.
\end{equation}
The constant coefficient of \(g/h\) at a zero \(\rho\) is
\[
\frac{g^{(m_\rho)}(\rho)}
 {m_\rho!\prod_{\sigma\ne\rho}(\rho-\sigma)^{m_\sigma}}\ne0 .
\]
Thus it is a unit in each truncated local ring: its
inverse is a finite geometric series after extracting
that nonzero constant. The substitution kernel is
\eqref{NH2}, as already calculated, so \(\eta\) is injective.
On its image, \(\eta T(t)\eta^{-1}\) has exactly the
heat operator \(\eta F_s(T(t))\eta^{-1}\), by its
absolutely convergent series; adjoints are transported
by the displayed metric. All determinants, traces,
multiplicities, and error constants above therefore
hold on this same physical image. The inverse is only
on \(\eta(C)\); no map on the ambient complement is asserted.
The attained source lift remains
\[
\mathcal A\sum_{j=0}^L(Rv)_jQ_j/\sqrt{\omega_j},
\qquad
\mathcal A P=P(D_1+\cdots+D_k)F_h^{\otimes k},
\]
with \(Q_j\) evaluated at \(u=(S-c)/i\).
Its class is \(\eta v\), by \(ER=I\), and its norm
squared is \(v^*Gv\). The heat calculation therefore
has not removed the original arithmetic source map.

\section{Exact finite comparisons, including defective endpoints}

The following matrices test the proved identities; they
are not presented as actual packets of zeta zeros.
For the separate comparison measure
\((3/2)1_{[-1,1]}(u)\,du\), of mass \(3\), and quotient
\(u^2+1\) at \(L=1\), retain
\begin{equation}\tag{NH27}\label{NH27}
G=\begin{pmatrix}3&0\\0&1\end{pmatrix},\quad
A=\begin{pmatrix}0&1/3\\1&0\end{pmatrix},\quad
Q=\begin{pmatrix}0&4/3\\0&0\end{pmatrix},\quad
a^2=1/3,\quad \eps^2=16/3,\quad \gamma_1=4/3 .
\end{equation}
Then \(T(t)^2=(1-4t)I/3\), so
\begin{equation}\tag{NH28}\label{NH28}
\Phi_s(t)=2e^{-s(1-4t)/3},\qquad
C_n(s)=2e^{-s/3}\frac{(4s/3)^n}{n!}.
\end{equation}
Here \(\gamma_{2j+1}=(4/3)(1/3)^j\).
There are \(\binom{n+r-1}{n-1}\) summands in \(H_{n,r}\);
inserting them into \eqref{NH15} sums to
\eqref{NH28}, including every sign.
At \(t=1/4\), \(T\ne0\), \(T^2=0\), so this is a
genuine size-two Jordan matrix although its heat trace
equals \(2\). The full operator is \(I\) here because
the even test has \(T^2=0\); the resolvent still has
its nonzero order-two term \(T/z^2\).
For the endpoint test quotient \(u^2\), replacing the
upper-right entry of \(Q\) by \(1/3\) gives
\begin{equation}\tag{NH29}\label{NH29}
T(1)=\begin{pmatrix}0&0\\1&0\end{pmatrix},\quad
\Phi_s(t)=2e^{-s(1-t)/3},\quad
\Phi_s^{(n)}(1)=2(s/3)^n,\quad
\kappa_1(z)=1/(3z^2).
\end{equation}
Direct substitution in \eqref{NH11} gives the same
derivative: \(F_s'(z)\kappa_1(z)^n\) has coefficient
\(2n(-s)^n/(n!\,3^n)\) on \(z^{-1}\), and multiplication
by \((-1)^n(n-1)!\) gives \(2(s/3)^n\).

For the four-dimensional comparison quotient
\((u^2+1)^2\), the complete retained matrices are
\begin{equation}\tag{NH30}\label{NH30}
\begin{gathered}
G=\begin{pmatrix}
3&0&1&0\\0&1&0&3/5\\1&0&3/5&0\\0&3/5&0&3/7
\end{pmatrix},\quad
M=\begin{pmatrix}
0&0&0&-1\\1&0&0&0\\0&1&0&-2\\0&0&1&0
\end{pmatrix},\\
r=(32/35,0,20/7,0)^{\mathsf T},\quad
\ell=(0,0,0,1),\quad Q=r\ell,\quad A=M+Q .
\end{gathered}
\end{equation}
Its entries \(G_{ij}=3/(i+j+1)\) for even \(i+j\)
and \(0\) otherwise are the actual moment integrals.
The polynomial
\(Q_4(u)=u^4-6u^2/7+3/35\) is orthogonal to
\(1,u,u^2,u^3\): the odd pairings vanish and the two
even pairings are
\(3/5-6/7+9/35=0\) and
\(3/7-18/35+3/35=0\).
Compression replaces \(u^4\) by \(6u^2/7-3/35\),
exactly the last column of \(A\), proving \(GA=A^*G\).
The original multiplication matrix is cyclic, with
minimal polynomial \((z^2+1)^2\). Moreover
\((M^2+I)^2=0\) while \(M^2+I\ne0\).
Thus both \(+i\) and \(-i\) have a size-two Jordan block.
Its full heat operator and trace are
\begin{equation}\tag{NH31}\label{NH31}
e^{-sM^2}=e^s\bigl(I-s(M^2+I)\bigr),\qquad
\Phi_s(1)=4e^s .
\end{equation}
This follows by inserting the square-zero \(M^2+I\)
in the exponential. Its trace is zero by nilpotence,
so the constant term has trace \(4e^s\), not \(2e^s\).
The full scalar endpoint remains
\[
\kappa_1(z)=
\frac{(20/7)z^2+32/35}{(z^2+1)^2},
\]
whose numerator at \(z=\pm i\) is \(-68/35\ne0\).
Thus the contour derivative formula retains both double poles.

\section{The metric heat difference recovers the same arithmetic allowance}

For real \(t\), keep the same adjoint and define
\begin{equation}\tag{NH32}\label{NH32}
\begin{gathered}
\Psi_s(t)=\tr e^{-sT(t)^\dd T(t)},\quad
\Delta_s(t)=\Re\Phi_s(t)-\Psi_s(t),\quad
b=\|T(t)\|_G,\quad X=sb^2,\\
\mathscr D(t)=\begin{pmatrix}0&T(t)\\T(t)^\dd&0\end{pmatrix}
\quad\hbox{on }(C\oplus C,G\oplus G),\\
\tr e^{-s\mathscr D(t)^2}=2\Psi_s(t),\qquad
P_1\mathscr D(t)I_2=T(t).
\end{gathered}
\end{equation}
The second line is a specified selfadjoint dilation, not
a replacement of \(\chi\). Its square has diagonal
blocks \(TT^\dd,T^\dd T\). Their traces of every positive
power agree by cyclicity, and their zero-th powers both
have trace \(q\). Their exponential traces agree by
absolute convergence, proving the factor \(2\).
The coordinate injection \(I_2v=(0,v)\) and projection
\(P_1(v,w)=v\) prove the last map directly.

\begin{theorem}[Finite, dimension-free recovery at explicit heat time]
The exact comparison satisfies
\begin{equation}\tag{NH33}\label{NH33}
\left|\Delta_s(t)-st^2\eps^2\right|
\le st^2\eps^2\bigl[(1+2X)e^X-1\bigr].
\end{equation}
In particular, for every \(s>0\) with \(sb^2\le1/7\),
\begin{equation}\tag{NH34}\label{NH34}
\frac12\,st^2\eps^2\le\Delta_s(t)
\le\frac32\,st^2\eps^2 .
\end{equation}
One may replace \(b\) in the sufficient heat-time bound
by the explicit upper bound \(a+|t|\eps\).
\end{theorem}
\begin{proof}
The complete independent calculation is also recorded
in \cite{MetricHeat}%
% BEGIN VERIFIED PUBLIC PROOF LINK
~\href{https://github.com/KokunoYumeto/zeta-function-research-reader/blob/5992ad50c0f2a06921f1fcaded7b4e38097c0739/workbenches/splitzero-tandem/continuations/20260920-original-kernel-mixed-determinants/providers/HEAT_METRIC_COMPARISON.tex\#L43}{[proof HM1]}%
% END VERIFIED PUBLIC PROOF LINK
{}; the required proof is supplied here.
Write \(T=H+iK\), with
\(H=(T+T^\dd)/2\) and \(K=(T-T^\dd)/(2i)\).
Both are selfadjoint. Since \(Q^2=(Q^\dd)^2=0\),
\(\tr K^2=t^2\eps^2/2\) by expanding its square.
For \(0\le v\le1\), the exact path
\[
Z(v)=H+ivK=\frac{1+v}{2}T+\frac{1-v}{2}T^\dd
\]
has norm at most \(b\), because \(\|T^\dd\|_G=\|T\|_G\).
Conjugation by \(G^{1/2}\) is an isometry to Euclidean
coordinates and sends the adjoint to the ordinary
conjugate transpose; it proves that norm equality
without removing \(G\).
Put
\[
f_n(v)=\tr(Z(v)^\dd Z(v))^n-\Re\tr Z(v)^{2n}.
\]
Then \(f_n(0)=0\). In the derivative of the first
trace at zero, the \(n\) positive and \(n\) negative
\(iK\) terms cancel by cyclicity. The derivative of the
second trace before taking the real part is
\(2ni\tr(H^{2n-1}K)\), whose real part vanishes because
the trace of the product of two selfadjoint matrices
is real. Thus \(f_n'(0)=0\).

Each second derivative of either length-\(2n\) word
has \(2n(2n-1)\) ordered terms. Each contains two
specified \(K\)'s; all its other factors have norm at
most \(b\). A cyclic word \(KUKV\) has absolute trace
at most \(\|K\|_{\mathrm{HS},G}^2\|U\|_G\|V\|_G\):
matrix-entry Cauchy--Schwarz gives
\(|\tr XY|\le\|X\|_{\mathrm{HS}}\|Y\|_{\mathrm{HS}}\),
and multiplication on either side bounds the
Hilbert--Schmidt norm by operator norm times
Hilbert--Schmidt norm. These statements follow in the
same isometric coordinates by summing squared column
norms. Integrate \(f_n''\) with weight \(1-v\),
whose integral is \(1/2\). For
\(D_n=\tr(T^\dd T)^n-\Re\tr T^{2n}=f_n(1)\), this gives
\[
|D_n|\le n(2n-1)t^2\eps^2 b^{2n-2}.
\]
Direct expansion at \(n=1\) gives the exact value
\(D_1=t^2\eps^2\).
The two absolutely convergent exponentials therefore give
\(\Delta_s=\sum_{n\ge1}(-1)^{n+1}s^nD_n/n!\).
The tail from \(n=2\) is bounded by
\[
st^2\eps^2
\sum_{n\ge2}\frac{n(2n-1)X^{n-1}}{n!}
=st^2\eps^2[(1+2X)e^X-1].
\]
The series identity is obtained by differentiating
\(e^X\) twice. If \(b=0\), then \(T=0\), hence
\(K=0\) and \(t^2\eps^2=0\); the identity holds without
dividing by \(b\). Finally \(e^X\le(1-X)^{-1}\) for
\(0\le X<1\), by its nonnegative series coefficients,
so \((1+2X)e^X-1\le3X/(1-X)\le1/2\) when
\(X\le1/7\). This proves \eqref{NH34}.
\end{proof}

The theorem has a direct original arithmetic receiver.
Let \(\lambda\) run over eigenvalues of \(B(t)=cI+iT(t)\)
with their full algebraic multiplicities, and put
\(\mathcal L(t)=\sum_{\Re\lambda>c}(2\Re\lambda-2c)\).
Then throughout the heat-time range in \eqref{NH34},
\begin{equation}\tag{NH35}\label{NH35}
\mathcal L(t)^2\le t^2\eps^2
\le\frac{2\Delta_s(t)}s
=\frac{2\Re\Phi_s(t)-\tr e^{-s\mathscr D(t)^2}}s .
\end{equation}
Here is the full finite spectral bound.
The Hermitian defect
\(B+B^\dd-2cI=it(Q^\dd-Q)\) has eigenvalues
\(t\eps,-t\eps\) and zero on the complement.
To see this explicitly, use the orthogonal columns
\(r/\|r\|_G\) and
\(G^{-1}\ell^*/\|G^{-1}\ell^*\|_G\), which are orthogonal
because \(\ell r=0\).
In that pair its matrix is
\(\left(\begin{smallmatrix}0&-it\eps\\it\eps&0\end{smallmatrix}\right)\).
If \(Q=0\) the whole defect is zero.
Let \(W\) be the sum of generalized eigenspaces of \(B\)
with real part above \(c\), and \(P_W\) its
\(G\)-orthogonal projection.
In a basis beginning with \(W\), invariance shows
\(\tr P_W(B+B^\dd-2cI)=\mathcal L(t)\).
Each nilpotent Jordan part has trace zero but retains
its full dimension in this identity.
In an orthonormal eigenbasis of the Hermitian defect,
each diagonal entry of \(P_W\) lies in \([0,1]\).
Its trace pairing is therefore at most the sum of
positive defect eigenvalues, namely \(|t|\eps\).
Squaring and using \eqref{NH34} proves \eqref{NH35}.

For any fixed nonzero \(t\), \eqref{NH33} also gives
the quantitative recovery
\(\eps^2=\lim_{s\downarrow0}\Delta_s(t)/(st^2)\).
This is a recovery of the existing arithmetic allowance
from two specified finite heat traces. It is not a
proof that the allowance decays in \(L\) or \(k\).
For the mass-three example \eqref{NH27}, the exact
observable, without any numerical approximation, is
\begin{equation}\tag{NH36}\label{NH36}
\Delta_s(t)=
2e^{-s(1-4t)/3}-e^{-s/3}
-e^{-s(1-4t)^2/3},\qquad
b^2=\max\{1/3,(1-4t)^2/3\}.
\end{equation}
Indeed \(T^\dd T=\operatorname{diag}(1/3,(1-4t)^2/3)\).
Its first heat coefficient is \(16t^2/3\), the
original \(t^2\eps^2\).

The holomorphic trace alone need not recover that norm.
The precise map is exhibited by
\(A=\operatorname{diag}(-2,-1,1,2)\), \(G=I_4\),
\[
r=(\alpha,\beta,\beta,\alpha)^{\mathsf T},\qquad
\ell=(\alpha^{-1},\beta^{-1},-\beta^{-1},-\alpha^{-1}),
\qquad \alpha,\beta>0.
\]
Its pairing \(\ell r\) is zero and its spectral weights
are always \(1,1,-1,-1\).
The diagonal conjugacy
\(D=\operatorname{diag}(\alpha,\beta,\beta,\alpha)\)
carries the case \(\alpha=\beta=1\) to these matrices,
commutes with \(A\), and carries every exponential by
the same conjugacy. Therefore \(\Phi_s(t)\) is identical
for all these choices, for every \(s,t\).
The squared fixed-metric norm, however, is
\(4(\alpha^2+\beta^2)(\alpha^{-2}+\beta^{-2})\):
it equals \(16\) at \((1,1)\) and \(25\) at \((2,1)\).
The conjugacy is not an isometry of the retained
\(I_4\), and \eqref{NH33} measures precisely this
change through the metric heat trace.
These are finite comparison families, not proposed
arithmetic packets.

\paragraph{What this calculation advances.}
The original outgoing relation now has a genuine entire
heat response with two explicit convergent expansions,
full original-coordinate and physical-image maps, and
factorial error bounds in the attained metric.
The same-metric heat difference now gives a quantitative
finite recovery of \(t^2\eps^2\), including the complete
exterior receiver \eqref{NH35}; it is not replaced by
\(\Phi_s\) alone. The independent companion
\href{https://github.com/KokunoYumeto/zeta-function-research-reader/blob/5992ad50c0f2a06921f1fcaded7b4e38097c0739/workbenches/splitzero-tandem/continuations/20260920-original-kernel-mixed-determinants/providers/HEAT_METRIC_COMPARISON.tex\#L1}{\texttt{HEAT\_METRIC\_COMPARISON.tex}} additionally retains
the next heat coefficient and its complete cubic term.
None of these finite formulas establishes decay of the
actual arithmetic \(\eps\) as \(L,k\) grow.

\begin{thebibliography}{9}\raggedright
\bibitem{ConnesHeat}
Alain Connes, \emph{Heat expansion and zeta},
\href{https://arxiv.org/abs/2402.13082v1}{arXiv:2402.13082v1}
(2024). Complete original author source
\texttt{heat-kernel-zeta.tex} read for this note.
Gaussian equation \texttt{ft}, heat identity \texttt{plicit1},
and the explicitly RH-assuming theorem
\texttt{heatkernelintro}; no infinite theorem is invoked.
Source SHA-256:
\texttt{e3d9bc846d966d19ceaf86133a863050c17c9bf40ba5b4a220945f0f5be89b42}.
\bibitem{CV}
Alain Connes and Walter D. van Suijlekom,
\emph{Quadratic Forms, Real Zeros and Echoes of the Spectral Action},
\href{https://arxiv.org/abs/2511.23257v1}{arXiv:2511.23257v1}
(2025), Section \texttt{sect:sa}, Lemma \texttt{lem:sa}.
This note uses the receiving derivation and source-class
check in \cite{NativeAction}; the complete original CV
source was read by that derivation's owner, not independently
read for this note. No source inspection beyond this is claimed.
\bibitem{NativeAction}
Split-Zero programme, \emph{Native spectral action along the
outgoing arithmetic relation}, 20 September 2026,
\href{https://github.com/KokunoYumeto/zeta-function-research-reader/blob/5992ad50c0f2a06921f1fcaded7b4e38097c0739/workbenches/splitzero-tandem/continuations/20260920-original-kernel-mixed-determinants/providers/NATIVE_SPECTRAL_ACTION.tex\#L1}{\texttt{NATIVE\_SPECTRAL\_ACTION.tex}}, SA1--SA43, with
\href{https://github.com/KokunoYumeto/zeta-function-research-reader/blob/5992ad50c0f2a06921f1fcaded7b4e38097c0739/workbenches/splitzero-tandem/continuations/20260920-original-kernel-mixed-determinants/providers/SOURCE_FUNCTION_CLASS_CHECK.tex\#L1}{\texttt{SOURCE\_FUNCTION\_CLASS\_CHECK.tex}}, FC1--FC4,
\href{https://github.com/KokunoYumeto/zeta-function-research-reader/blob/5992ad50c0f2a06921f1fcaded7b4e38097c0739/workbenches/splitzero-tandem/continuations/20260920-original-kernel-mixed-determinants/providers/RANK_ONE_METRIC_CURVATURE.tex\#L1}{\texttt{RANK\_ONE\_METRIC\_CURVATURE.tex}}, MC1--MC35, and
\href{https://github.com/KokunoYumeto/zeta-function-research-reader/blob/5992ad50c0f2a06921f1fcaded7b4e38097c0739/workbenches/splitzero-tandem/continuations/20260920-original-kernel-mixed-determinants/providers/INDEPENDENT_ENDPOINT_CHECK.tex\#L1}{\texttt{INDEPENDENT\_ENDPOINT\_CHECK.tex}}, EC1--EC15.
These complete programme proofs were read directly.
They are local programme derivations, not independent
human scholarly authorities.

% BEGIN VERIFIED PUBLIC PROOF LINK
\href{https://github.com/KokunoYumeto/zeta-function-research-reader/blob/5992ad50c0f2a06921f1fcaded7b4e38097c0739/workbenches/splitzero-tandem/continuations/20260920-original-kernel-mixed-determinants/providers/NATIVE_SPECTRAL_ACTION.tex\#L31}{Source proof SA1}. \href{https://github.com/KokunoYumeto/zeta-function-research-reader/blob/5992ad50c0f2a06921f1fcaded7b4e38097c0739/workbenches/splitzero-tandem/continuations/20260920-original-kernel-mixed-determinants/providers/SOURCE_FUNCTION_CLASS_CHECK.tex\#L24}{Source proof FC1}. \href{https://github.com/KokunoYumeto/zeta-function-research-reader/blob/5992ad50c0f2a06921f1fcaded7b4e38097c0739/workbenches/splitzero-tandem/continuations/20260920-original-kernel-mixed-determinants/providers/RANK_ONE_METRIC_CURVATURE.tex\#L35}{Source proof MC1}. \href{https://github.com/KokunoYumeto/zeta-function-research-reader/blob/5992ad50c0f2a06921f1fcaded7b4e38097c0739/workbenches/splitzero-tandem/continuations/20260920-original-kernel-mixed-determinants/providers/INDEPENDENT_ENDPOINT_CHECK.tex\#L28}{Source proof EC1}.
% END VERIFIED PUBLIC PROOF LINK
\bibitem{MetricHeat}
Split-Zero programme, \emph{Recovering the original rank-one
metric allowance from two finite heat traces},
20 September 2026,
\href{https://github.com/KokunoYumeto/zeta-function-research-reader/blob/5992ad50c0f2a06921f1fcaded7b4e38097c0739/workbenches/splitzero-tandem/continuations/20260920-original-kernel-mixed-determinants/providers/HEAT_METRIC_COMPARISON.tex\#L1}{\texttt{metric\_check/HEAT\_METRIC\_COMPARISON.tex}},
HM1--HM19. Independent finite derivation read in full
for the comparison proved again in \eqref{NH32}--\eqref{NH36}.
It is a programme derivation, not a human primary source.

% BEGIN VERIFIED PUBLIC PROOF LINK
\href{https://github.com/KokunoYumeto/zeta-function-research-reader/blob/5992ad50c0f2a06921f1fcaded7b4e38097c0739/workbenches/splitzero-tandem/continuations/20260920-original-kernel-mixed-determinants/providers/HEAT_METRIC_COMPARISON.tex\#L43}{Source proof HM1}.
% END VERIFIED PUBLIC PROOF LINK
\end{thebibliography}
\end{document}
